% =============================================================================
%  Notas del Profesor — Sesión 01: Presentación del Curso
%  Programación Avanzada — Universidad Panamericana
%  Compilar: pdflatex sesion01_presentacion_profesor.tex
% =============================================================================
\documentclass[12pt, oneside]{article}
\usepackage[profesor]{progavanzada}

\begin{document}

\portada{Sesión 01}{Presentación del Curso}{2026}

\tableofcontents
\newpage

% =============================================================================
\section{Objetivos de la sesión}
% =============================================================================

\begin{enumerate}
  \item Presentar al profesor, el curso y sus expectativas.
  \item Transmitir el syllabus, la bibliografía y la forma de calificación.
  \item Conocer al grupo: niveles de experiencia, expectativas y emociones
        iniciales frente a la programación.
  \item Establecer desde el primer día una cultura de participación activa.
  \item Medir la velocidad de escritura como diagnóstico inicial (10FastFingers).
\end{enumerate}

\begin{notaprofesor}[Tono de la sesión]
  Esta clase es más de clima y contrato pedagógico que de contenido técnico.
  El objetivo principal es que los alumnos salgan con ganas de volver.
  Sé accesible, muestra entusiasmo genuino por el tema y valida que la
  programación puede ser difícil — y que eso está bien. Evita el discurso
  de ``este curso es muy difícil, muchos van a reprobar'': genera ansiedad
  sin ningún beneficio pedagógico.
\end{notaprofesor}

% =============================================================================
\section{Duración y materiales}
% =============================================================================

\begin{center}
\begin{tabular}{lc}
  \toprule
  \textbf{Actividad} & \textbf{Tiempo estimado} \\
  \midrule
  Presentación del profesor y el curso & 10 min \\
  Syllabus y calificación              & 15 min \\
  Bibliografía y recursos              & 10 min \\
  Padlet — meme de programar           & 15 min \\
  Lluvia de ideas: ¿buen programador?  & 20 min \\
  10FastFingers                        & 10 min \\
  Cierre y vista al siguiente clase    &  5 min \\
  \midrule
  \textbf{Total}                       & \textbf{85 min} \\
  \bottomrule
\end{tabular}
\end{center}

\noindent\textbf{Materiales necesarios:}
\begin{itemize}
  \item Proyector con acceso a internet
  \item Padlet creado antes de la clase (ver instrucciones en \S3)
  \item Los alumnos con laptop, tablet o smartphone para acceder al Padlet
        y a 10FastFingers
\end{itemize}

% =============================================================================
\section{Secuencia detallada}
% =============================================================================

\subsection{Presentación del profesor y del curso (10 min)}

Comienza con una breve presentación personal: quién eres, en qué trabajas o
has trabajado, qué te apasiona de C++ y la programación. No más de 3 minutos.

Luego una pregunta de sondeo rápida al grupo (manos levantadas):
\begin{itemize}
  \item ¿Cuántos han programado antes en algún lenguaje?
  \item ¿Cuántos han visto C o C++ anteriormente?
  \item ¿Cuántos usan regularmente herramientas de IA para programar?
\end{itemize}

\begin{notaprofesor}
  El sondeo de IA es especialmente revelador en grupos actuales. Si la
  mayoría usa Copilot o ChatGPT, ajusta el discurso del curso: enfatiza
  que entender la memoria y los apuntadores es exactamente lo que distingue
  a quien puede \textit{evaluar} el código que genera la IA del que solo
  puede copiarlo.
\end{notaprofesor}

\subsection{Syllabus y calificación (15 min)}

Presenta el syllabus proyectado. Detente especialmente en la estructura de
calificación: 1/3 tareas, 1/3 asistencia y participación, 1/3 examen.

\begin{notaprofesor}[Por qué este esquema de calificación]
  El 33\% de asistencia y participación no es arbitrario. La programación
  se aprende haciendo y discutiendo, no leyendo. Un alumno presente y activo
  que no llega al 70 en el examen tiene más probabilidad de resolver problemas
  reales que uno que memorizó para el examen pero nunca participó. Explica
  esto explícitamente: muchos alumnos no están acostumbrados a que la
  participación valga tanto.

  \vspace{0.5em}
  Define claramente qué es ``participación'': hacer preguntas, responder,
  proponer ideas, ayudar a un compañero frente al grupo. No es asistencia
  silenciosa.
\end{notaprofesor}

Abre espacio para preguntas sobre el reglamento. Es mejor aclarar dudas
sobre entregas, fechas y política de retrasos en la primera clase.

\subsection{Bibliografía y recursos (10 min)}

Proyecta cppreference.com y learncpp.com. Muestra rápidamente cómo están
organizados. No es necesario entrar en detalle — la intención es que sepan
que existen y que los usarán durante el semestre.

\begin{notaprofesor}
  Muchos alumnos jamás han consultado documentación técnica formal. Decirles
  ``lean el libro'' sin enseñarles a navegarlo es ineficaz. En la sesión 2
  harás una demostración en vivo de cppreference — aquí solo siémbra la idea.
\end{notaprofesor}

\subsection{Padlet — ¿Qué emoción te produce programar? (15 min)}

\textbf{Configuración previa (hacer antes de la clase):}
\begin{enumerate}
  \item Crear cuenta en \url{https://padlet.com} si no tienes una.
  \item Nuevo padlet $\to$ formato \textbf{Muro} (Wall).
  \item Título: ``¿Qué emoción te produce programar?''
  \item Activar opción de comentarios si deseas discusión.
  \item Copiar el enlace o generar el QR code (botón Share).
  \item Proyectar el QR al inicio de la actividad.
\end{enumerate}

\textbf{Instrucción al grupo:}

\begin{quote}
  ``Busca un meme, GIF o imagen que represente cómo te sientes cuando
  programas — o cómo imaginas que te vas a sentir. Agrégalo al Padlet
  con tu nombre. Tienes 8 minutos.''
\end{quote}

Dedica los últimos 7 minutos a revisar los posts con todo el grupo. Comenta
algunos, genera risa, normaliza la frustración y el caos que aparezcan.

\begin{notaprofesor}[Para qué sirve esta actividad]
  Tiene tres funciones simultáneas: rompe el hielo, da información diagnóstica
  sobre la actitud del grupo frente a la programación, y establece que este
  curso va a tener un tono humano. Los alumnos que suben memes de frustración
  total son frecuentemente los más comprometidos al final del semestre — solo
  necesitaban permiso para expresarlo.

  \vspace{0.5em}
  Guarda el Padlet. Es un documento vivo: puedes mostrarlos al final del
  semestre y pedir a los alumnos que comenten si su meme inicial todavía
  los representa.
\end{notaprofesor}

\subsection{Lluvia de ideas: ¿Qué se necesita para ser un buen programador hoy? (20 min)}

Escribe la pregunta en el pizarrón tal cual. Sin definirla, sin restringirla.

\textbf{Dinámica:}
\begin{enumerate}
  \item 3 minutos de reflexión individual en silencio.
  \item 10 minutos de lluvia de ideas abierta: cualquiera habla, tú escribes
        en el pizarrón sin filtrar ni jerarquizar.
  \item 7 minutos de discusión: agrupa las ideas que aparezcan en categorías
        naturales (técnicas, actitudinales, sociales, etc.).
\end{enumerate}

\begin{notaprofesor}[Ideas que suelen emerger — y cómo guiarlas]

  \textbf{``Saber muchos lenguajes''} — Guía hacia: la profundidad en uno
  aporta más que la superficialidad en muchos. C++ enseña conceptos que
  aplican a todos los lenguajes.

  \vspace{0.4em}
  \textbf{``Usar IA / Copilot''} — Valida: es una herramienta real. Agrega:
  para evaluar lo que genera la IA necesitas entender qué está haciendo. Un
  programador que no entiende apuntadores no puede juzgar si el código
  generado gestiona bien la memoria.

  \vspace{0.4em}
  \textbf{``Paciencia / no rendirse''} — Valida completamente. Es real y
  está respaldado por la literatura: la perseverancia predice el éxito en
  programación más que la aptitud inicial.

  \vspace{0.4em}
  \textbf{``Matemáticas''} — Matiza: para programación general, no es un
  requisito estricto. Para ciertos campos (gráficos, ML, criptografía, física)
  sí importa más. En este curso las matemáticas que necesitarás son mínimas.

  \vspace{0.4em}
  \textbf{``Buscar en Google / Stack Overflow''} — Valida: es una habilidad
  real. Saber formular una búsqueda efectiva es parte del oficio.
\end{notaprofesor}

\begin{notaprofesor}[Dato para compartir con el grupo]
  Un estudio de la Universidad de Washington (2019) siguió a estudiantes de
  CS durante 4 años y encontró que la variable más predictiva de completar la
  carrera no fue la nota de entrada ni el conocimiento previo, sino la
  \textbf{mentalidad de crecimiento} (growth mindset): creer que la
  inteligencia y la habilidad se desarrollan con esfuerzo.\\
  \textit{Referencia:} Cutts et al. (2019). ``Predict or Pre-empt: How to
  respond to poor performance in early CS courses''. \textit{ICER 2019}.
\end{notaprofesor}

\subsection{10FastFingers — Diagnóstico de mecanografía (10 min)}

URL: \url{https://10fastfingers.com/typing-test/spanish}

\textbf{Instrucción al grupo:}
\begin{quote}
  ``Abre el sitio, elige el modo en español y haz el test de 1 minuto.
  Anota tus palabras por minuto y tu precisión en tu hoja de notas.
  Nadie tiene que mostrar su resultado si no quiere.''
\end{quote}

\begin{notaprofesor}[Por qué medir mecanografía]
  No es un criterio de calificación. Su valor es diagnóstico y motivacional:

  \begin{itemize}
    \item Diagnostica cuánto tiempo perderá el alumno en actividades simples
          de escritura durante el semestre.
    \item Abre la conversación sobre herramientas y hábitos: atajos de teclado,
          teclado mecánico, práctica deliberada.
    \item Permite una comparación al final del semestre: los alumnos suelen
          mejorar 15--25 PPM solo por la práctica acumulada de escribir código.
  \end{itemize}

  Si quieres hacer seguimiento, pide voluntarios que compartan su PPM inicial
  en el Padlet. Puedes repetir el ejercicio al final del semestre.
\end{notaprofesor}

Comparte con el grupo algunos datos de contexto (ver notas del alumno,
\S4.3). La idea no es presionar sino motivar: la mecanografía mejora
dramáticamente con práctica deliberada de 10--15 minutos al día.

\subsection{Cierre y vista a la siguiente clase (5 min)}

\begin{itemize}
  \item Resume en una frase qué viene en la sesión 2: ``La próxima clase
        escribimos y ejecutamos nuestro primer programa C++.''
  \item Recuerda el requisito técnico: \texttt{g++} instalado o acceso a
        Google Cloud Shell.
  \item Deja abierta la pregunta de la lluvia de ideas: ``Vayan pensando en
        si añadirían algo a la lista que armamos hoy.''
\end{itemize}

% =============================================================================
\section{Adaptaciones según el perfil del grupo}
% =============================================================================

\begin{notaprofesor}[Grupo sin experiencia previa]
  Dedica más tiempo al Padlet y a la lluvia de ideas. La ansiedad ante lo
  desconocido es real; nombrarla y normalizarla es más efectivo que minimizarla.
  Puedes agregar una ronda de presentaciones personales al inicio.
\end{notaprofesor}

\begin{notaprofesor}[Grupo con experiencia en otros lenguajes]
  El sondeo inicial revelará esto. Ajusta el discurso: C++ no es ``otro
  lenguaje más'' — es un cambio de modelo mental. La gestión explícita de
  memoria y la falta de garbage collector son el salto conceptual más
  importante. Puedes hacer la pregunta de lluvia de ideas más específica:
  ``¿Qué se necesita para \textit{dominar} C++ viniendo de Python/JavaScript?''
\end{notaprofesor}

% =============================================================================
\section{Recursos adicionales para el profesor}
% =============================================================================

\begin{itemize}
  \item \textbf{Padlet (guía rápida):} \url{https://padlet.com/support}
  \item \textbf{10FastFingers:} \url{https://10fastfingers.com}
  \item \textbf{Índice TIOBE C++:} \url{https://www.tiobe.com/tiobe-index/cplusplus/}
  \item \textbf{Growth mindset en CS:} Dweck, C. (2006). \textit{Mindset: The New
        Psychology of Success}. Random House.
  \item \textbf{C++ Core Guidelines (referencia de calidad de código):}\\
        \url{https://isocpp.github.io/CppCoreGuidelines/CppCoreGuidelines}
\end{itemize}

% =============================================================================
\section{Material de la sesión}
% =============================================================================

\begin{itemize}
  \item \textbf{Notas del alumno:}
        \texttt{notas\_alumno/pdf/sesion01\_presentacion.pdf}
  \item \textbf{Siguiente sesión:}
        \texttt{sesiones/sesion02\_introduccion\_cpp.ipynb}
\end{itemize}

\end{document}
